\begin{abstract}
Time-series foundation models now forecast any series zero-shot, while graph specialists are still trained per
network, and the two are rarely compared under one protocol. We compare them on physical sensor networks with a fixed,
open benchmark: six traffic networks (PEMS03/04/07/08, METR-LA, PEMS-BAY) and two monsoon rainfall regions, four
tracks (standard context, open context, 30\% dead sensors, zero-shot transfer), one scorer for point error, pinball
loss and interval score, saved forecasts and paired Diebold--Mariano and block-bootstrap tests. We introduce PWM-Dyn,
a 0.3M-parameter graph-diffusion model with explicit missing-data inputs, quantile heads and split-conformal
calibration, and compare it with STID, Graph WaveNet and STAEformer trained with their official code on identical
windows under a fixed compute budget, and with zero-shot TimesFM 2.5 and 3.0. Configurations were chosen on
validation and the reported models use fresh seeds; four networks were never used during development. Under a
12-step context PWM-Dyn has the lowest MAE on five of six traffic networks (0.3--4.2\% below the best trained opponent
on the four new ones), is statistically tied with STID on two, and trails a budget-capped STAEformer on PEMS08. Against
TimesFM 3.0 it has lower MAE and interval score on seven of eight tasks and loses on PEMS03. With 30\% of sensors dark it beats per-sensor
climatology on seven of eight tasks, where specialists without a missing-data mechanism fail. An ablation traces an earlier
accuracy loss to a climatology input. Code, benchmark and forecasts are released.
\end{abstract}
